\documentclass{article}
\usepackage{iclr2025_conference,times}
\usepackage{hyperref}
\usepackage{url}
\usepackage{graphicx}
\usepackage{array}
\usepackage{amsmath}
\usepackage{multirow}
\usepackage{multicol}
\usepackage{booktabs}
\usepackage{tabularx}
\usepackage{longtable}
\usepackage{listings}
\usepackage{tcolorbox}
\usepackage{longtable}
\usepackage[utf8]{inputenc}
\usepackage{lipsum} % For dummy text
\usepackage{subcaption}
\usepackage[normalem]{ulem}
\usepackage{enumitem}
\usepackage{adjustbox}
\usepackage{wrapfig,lipsum,booktabs}
\usepackage{setspace}
\usepackage{tikz}
\usepackage{colortbl}
\usepackage{authblk}
\usepackage{CJKutf8}
\usepackage[T1]{fontenc}


\iclrfinaltrue
\renewcommand\Authfont{\bfseries} % 使作者名字为粗体

\renewcommand\Authfont{\bfseries} % 使作者名字为粗体
\title{Chain of Ideas: Revolutionizing Research via Novel Idea Development with LLM Agents}

\author{
Long Li\raisebox{4pt}{\small $1,2,{\thanks{This work was done when Long Li was an intern at Alibaba DAMO Academy}}$} \quad
Weiwen Xu\raisebox{4pt}{\small $1$} \quad
Jiayan Guo\raisebox{4pt}{\small $1$} \quad
Ruochen Zhao\raisebox{4pt}{\small $1$} \quad
Xingxuan Li\raisebox{4pt}{\small $1$} \quad
Yuqian Yuan\raisebox{4pt}{\small $1,2$} \quad
Boqiang Zhang\raisebox{4pt}{\small $1,3$} \quad
Yuming Jiang\raisebox{4pt}{\small $1$} \quad
Yifei Xin\raisebox{4pt}{\small $1$} \quad
Ronghao Dang\raisebox{4pt}{\small $1$} \quad
Yu Rong\raisebox{4pt}{\small $1$} \quad
Deli Zhao\raisebox{4pt}{\small $1$} \quad
Tian Feng\raisebox{4pt}{\small $2$} \quad
\textbf{Lidong Bing}\raisebox{4pt}{\small $1,\thanks{Lidong Bing is the corresponding author}$} \\
\raisebox{4pt}{\small $1$}DAMO Academy, Alibaba Group \\
\raisebox{4pt}{\small $2$}Zhejiang University \\
\raisebox{4pt}{\small $3$}University of Science and Technology of China \\
{\tt longli@zju.edu.cn} \\
{\tt weiwen.xuu@gmail.com} \\
{\tt binglidong@gmail.com}
}


\date{September 2024}

% \author{
% Long Li\raisebox{4pt}{\small $12^{\dagger}$} Weiwen Xu\raisebox{4pt}{\small $12^{\dagger}$} \quad Xin Li\raisebox{4pt}{\small $2^\ddagger$}  \quad Wenxuan Zhang\raisebox{4pt}{\small $2$} \quad  Meng Zhou\raisebox{4pt}{\small $3^{\dagger}$}\\ 
% \textbf{Wai Lam}\raisebox{4pt}{\small $1$} \quad  \textbf{Luo Si}\raisebox{4pt}{\small $2$}  \quad \textbf{Lidong Bing}\raisebox{4pt}{\small $2$} \\
% \raisebox{4pt}{\small $1$}DAMO Academy, Alibaba Group \\
% \raisebox{4pt}{\small $2$}Zhejiang University \\
% \raisebox{4pt}{\small $3$}University of Science and Technology of China \\
% {\tt longli@zju.edu.cn}\\
% {\tt weiwen.xuu@gmail.com}\\
% {\tt binglidong@gmail.com}
% }


\newcommand{\orange}[1]{\textcolor{orange}{#1}}
\newcommand{\red}[1]{\textcolor{red}{#1}}
\newcommand{\bing}[1]{\textcolor{brown}{[Bing: #1]}}
\newcommand{\xww}[1]{\textcolor{blue}{[xww: #1]}}
\newcommand{\gjy}[1]{\textcolor{green}{[gjy: #1]}}
\newcommand{\zrc}[1]{\textcolor{magenta}{[zrc: #1]}}
\newcommand*{\mybox}[2]{\tikz[anchor=base,baseline=0pt,rounded corners=0pt, inner sep=0.2mm] \node[fill=#1] (X) {#2};}

\newcommand{\cchar}[1]{\begin{CJK*}{UTF8}{gkai}#1\end{CJK*}}
\newcommand{\bingc}[1]{\textcolor{brown}{[\textbf{BING:} \cchar{#1}]}}

\begin{document}

\maketitle

\begin{abstract}

Effective research ideation is a critical step for scientific research. However, the exponential increase in scientific literature makes it challenging for researchers to stay current with recent advances and identify meaningful research directions. Recent developments in large language models~(LLMs) suggest a promising avenue for automating the generation of novel research ideas. However, existing methods for idea generation either trivially prompt LLMs or directly expose LLMs to extensive literature without indicating useful information. 
Inspired by the research process of human researchers, we propose a Chain-of-Ideas~(CoI) agent, an LLM-based agent that organizes relevant literature in a chain structure to effectively mirror the progressive development in a research domain.
This organization facilitates LLMs to capture the current advancements in research, thereby enhancing their ideation capabilities.
Furthermore, we propose Idea Arena, an evaluation protocol that can comprehensively evaluate idea generation methods from different perspectives, aligning closely with the preferences of human researchers.
Experimental results indicate that the CoI agent consistently outperforms other methods and shows comparable quality as humans in research idea generation. 
Moreover, our CoI agent is budget-friendly, with a minimum cost of \$0.50 to generate a candidate idea and its corresponding experimental design\footnote{The code is released at \url{https://github.com/DAMO-NLP-SG/CoI-Agent}. Try our online demo at \url{https://huggingface.co/spaces/DAMO-NLP-SG/CoI_Agent}.}.

\end{abstract}
% \gjy{ready for review.}

% \gjy{\cchar{感觉可以把arena-style evluation起名叫Idea Arena，再强调一下贡献}} \bing{\cchar{有道理！}}




% \bing{use CoI Agent to mention our framework.}
% \xww{Ready for review}\bing{did one pass}

% \orange{Why researchers believe LLMs could be a good tool for helping researchers conduct research}
\section{Introduction}
% \bingc{reminders:
% \begin{itemize}
%     \item In the appendix, the tables should be ordered according to the reference order in the main paper.
% \end{itemize}
% }

Idea generation is a crucial aspect of scientific research for driving technological innovations and breakthroughs. Traditionally, this process has been predominantly human-driven, necessitating expert researchers to review extensive literature, identify limitations in existing solutions, and propose new research directions. However, the complexity and vastness of scientific literature, coupled with rapid technological advancements, have rendered this task increasingly challenging for researchers.

% Idea generation is key to driving scientific innovation. Traditionally, this has relied on expert researchers to review vast literature, identify gaps, and suggest new directions. However, with the growing complexity of research and rapid technological advances, this task has become increasingly difficult.

Recent advancements in large language models~(LLMs) \citep{achiam2023gpt4,dubey2024llama3,yang2024qwen2} have enabled these models to exceed human experts in various scientific tasks, including mathematics \citep{yu2023metamath}, theorem proving \citep{yang2023leandojo}, and coding \citep{chen2021evaluating}. Building on this robust scientific foundation, one may hypothesize that LLMs could support a more abstract and creative research idea-generation task. Notably, \cite{si2024stanfordcan, kumar2024can} have validated this hypothesis, highlighting its substantial potential to expedite the discovery of novel concepts and uncharted research avenues.

%Recently, the rapid evolution in large language models' (LLMs) capabilities \citep{achiam2023gpt4,dubey2024llama3,Anthropic2024claude3,yang2024qwen2} has enabled LLMs to surpass human experts in various scientific tasks, including mathematics, theorem proving, and coding. Building upon this strong scientific foundation, it is reasonable to hypothesize that LLMs could facilitate more abstract and creative scientific tasks, such as idea generation. Encouragingly,  \citep{si2024stanfordcan, kumar2024can} have proven that LLMs can generate research ideas that closely resemble human creativity. This advancement has the potential to accelerate the discovery of novel concepts and unexplored research directions.

% \orange{Existing methods and their limitations. A key limitation should be that they cannot conduct idea derivation in the way that we humans do.}

Existing methods seek to address two key challenges to improve the quality of generated ideas: curating pertinent literature for LLMs to gain inspiration and ensuring the novelty of generated ideas. To address the first challenge, previous research enhances traditional academic retrieval systems, which typically depend on textual similarity, with academic knowledge graphs \citep{baek2024researchagent,wang2023scimon}. 
For the second challenge, existing approaches either apply predefined criteria such as novelty to guide the idea generation process \citep{baek2024researchagent} or iteratively refine ideas until they demonstrate low embedding similarities with existing papers \citep{wang2023scimon}.


% \gjy{In previous methods, LLMs are merely exposed to a large volume of information to generate the idea, where and the process of how to generate ideas are totally designed }
However, in existing attempts, LLMs are presented with an extensive volume of research literature when asked to generate ideas. This makes LLMs vulnerable to the influence of less relevant works, potentially resulting in ideas that lack logical coherence and technological innovation.
As shown in the upper part of Figure~\ref{fig:intro}, the LLM borrows an idea from GraphGPT \citep{tang2024graphgpt} and applies it into GoT framework~\citep{besta2024graph} to generate what they interpret as a ``novel idea''. However, the resultant idea conflates two concepts: GoT is a prompting method while GraphGPT is a fine-tuning method leveraging graph neural network architecture \citep{zhou2020graph}.
In contrast, human researchers often trace the evolution of a research field by analyzing its progression from foundational works to the most recent advancements. This comprehensive perspective provides valuable insights into the key factors driving developments within the domain. Such an understanding enables researchers to critically assess the limitations of earlier studies while identifying emerging trends. Therefore, they are better grounded in devising innovative and impactful research ideas.
\begin{figure}[t] % 设定图片的浮动位置
    \centering % 使图片居中
    \includegraphics[width=0.95\textwidth]{chaps/imgs/intro.pdf} % 插入图片，width控制宽度
    \caption{Comparison between the vanilla retrieval augmented generation (RAG) research agent and our Chain-of-Ideas agent on the idea generation task. } % 图片的标题
    \label{fig:intro} % 图片的标签，用于引用
    \vspace{-15pt}
\end{figure}

% \orange{Why Chain of Ideas is a good way to overcome those limitations of existing methods: Chain of Ideas is natural and exactly the way we human researchers do; LLMs are good at making step-by-step reasoning (e.g. CoT) for logically coherent derivation.}
% The process by which researchers develop new ideas through a comprehensive analysis of existing research threads underscores the progressive nature of human cognition. 
Motivated by the human practices in conducting research, we introduce a novel Chain-of-Ideas (CoI) agent framework to address the previously identified logical inconsistencies in the ideation processes of LLMs.
As shown in the bottom part of Figure \ref{fig:intro}, CoI agent aims to provide a clear landscape of current research topics by systematically selecting and organizing the relevant papers and their ideas in a chain structure.
CoI agent offers several distinctive advantages:
Firstly, it minimizes the risk of interference from less relevant literature via carefully selecting papers (i.e. from CoT \citep{wei2022chain} to GoT).
Second, LLMs are demonstrated with human practice to craft a novel idea. For example, SC \citep{wang2022self} emerges as a novel idea derived from CoT. This can be viewed as a form of few-shot prompting strategy, which has been proven to enhance the overall LLM's generation capability \citep{brown2020languagefewshot}. 
Third, CoI exemplifies a global progression in research development. As a result, LLMs can gain a deep understanding of the motivations behind these developmental trends, facilitating the identification of promising future research directions.

% \bing{The process by which researchers develop new ideas through a comprehensive analysis of existing research threads closely resembles the Chain of Thought (CoT) methodology, a well-established approach that has demonstrated effectiveness in guiding large language models (LLMs) in tackling complex tasks. In this paper, we introduce the term Chain of Ideas (CoI) to highlight the evolutionary nature of research ideas. By utilizing CoI to guide the step-by-step reasoning of LLMs in generating research ideas, we can enhance the coherence, relevance, and innovativeness of their outputs. By organizing research development in a chain structure that reflects the progressive trajectory, CoI facilitates a clearer understanding of research trends and promotes logical reasoning regarding the research evolution. Therefore, this approach mirrors the cognitive process that researchers naturally employ when formulating new concepts. Furthermore, it addresses the limitations of existing methodologies by alleviating the burden on LLMs that arise from direct exposure to vast amounts of information when generating new ideas.}
% This process by which human researchers generate new research ideas by analyzing existing ones step-by-step closely reflects chain of thought~(CoT), a well known approach that has proven effective in guiding LLMs when handling complex tasks. We call this process chain of idea~(CoI). CoI effectively leverages the strength that LLM excels at step-by-step reasoning. Prompting LLMs through CoI supports LLMs in producing more coherent, relevant, and innovative outputs by simulating the natural cognitive flow human researchers follow when developing new ideas. By organizing information in a chain structure that reflects the progressive relationships between ideas, CoI enables a clear understanding of research trends and supports logical reasoning about idea evolution. It also overcome the limitation of existing approaches by reliving the burden of LLMs that are directly exposing to a large volume of information to generate new ideas.

% \orange{Our framework design}


Specifically, CoI agent first retrieves an anchor paper of the given research topic. Instead of indiscriminately aggregating all papers within the citation network of the anchor, as done in \citep{baek2024researchagent}, we construct the CoI by selecting relevant and important literature from both the anchor's references and its subsequent works, thereby extending the chain backward and forward from the anchor.
We then prompt the constructed CoI to an LLM for idea generation and experiment design.
During idea generation, we require the LLM to predict possible future trends.
This prognostic result facilitates the gradual consolidation of the idea, beginning with the motivation for the proposed idea, progressing through an assessment of its potential impact, and culminating in the realization.
However, as the evolution of scientific discovery can emerge from multiple perspectives, a single CoI may be insufficient to capture the most promising direction. Additionally, there is no guarantee that the generated ideas will be novel. 
To address these issues, we construct multiple CoI branches for different perspectives of a research topic.
Additionally, a novelty-checker agent iteratively evaluates the draft idea against existing literature and refines it if substantial similarity is identified.

% \orange{Our main results and key findings} ~ \\

We compare our CoI agent against existing baselines on idea generation in the artificial intelligence (AI) field. To do this, we develop an arena-style evaluation framework called Idea Arena where participant methods compete in pairs, which demonstrates high agreement with human evaluation. The experimental results show that CoI agent consistently ranks first among all automated baselines, surpassing the second-best one by 56 ELO scores in human evaluation. CoI agent can generate ideas as novel as those of human experts. 
Our analysis further shows that for LLMs to generate novel ideas, a clear developmental trend analysis is more pivotal than the quantity of related literature.

Our contributions are summarized as follows: 
1) We propose the CoI agent to enhance LLMs' capability in idea generation. CoI agent  organizes relevant literature in a chain structure to effectively mirror the progressive nature of research development, allowing LLMs to better grasp the current  research advancements. 2) We propose Idea Arena for a comprehensive evaluation of idea-generation methods, which shows high agreement with human researchers. 3) Extensive experiments demonstrate the effectiveness of our CoI agent in generating ideas that are comparable to human creativity.


% \input{chaps/introduction}



\section{\textsc{SciExplore}}
% To operationalize the desiderata above, \textsc{SciExplore} instantiates scientific information-seeking as four complementary tasks. We discuss the task definitions in Section~\ref{sec:task_definition}, data construction in Section~\ref{sec:dataset_construction}, and quality control in Section~\ref{sec:quality_control}. Detailed data statistics and evaluation metrics can be found in Appendix~\ref{appendix_dataset_distribution} and Appendix~\ref{app:evaluation_details}, respectively.

To operationalize the desiderata above, \textsc{SciExplore} instantiates scientific information-seeking as four complementary tasks, curated through a fully expert-driven construction and verification pipeline. We formalize this task taxonomy in Section~\ref{sec:task_definition}, detail the corresponding data construction procedure in Section~\ref{sec:dataset_construction}, and describe our quality control mechanisms in Section~\ref{sec:quality_control}. Detailed data statistics and evaluation metrics are further provided in Appendix~\ref{appendix_dataset_distribution} and Appendix~\ref{app:evaluation_details}, respectively.

% To operationalize the desiderata above, \textsc{SciExplore} instantiates scientific information-seeking as four complementary tasks---spanning database navigation, ambiguous literature retrieval, missing reference completion, and cross-source structured knowledge synthesis---that jointly stress an agent's ability to locate, disambiguate, and integrate evidence from authentic scientific sources. Every instance is authored and verified by domain experts, and the full pipeline is designed to enforce realism, answer uniqueness, and resistance to parametric shortcuts, so that strong performance necessarily reflects genuine retrieval and reasoning rather than memorization. 
% We discuss the task definitions in Section~\ref{sec:task_definition}, data construction in Section~\ref{sec:dataset_construction}, and quality control in Section~\ref{sec:quality_control}. Detailed data statistics and evaluation metrics can be found in Appendix~\ref{appendix_dataset_distribution} and Appendix~\ref{app:evaluation_details}, respectively.

% \subsection{Task Definition and Construction}

\subsection{Task Definition}\label{sec:task_definition}


% Motivated by the limitations of existing benchmarks discussed above, \textsc{SciExplore} targets scientific information-seeking behaviors that arise in authentic research workflows but are largely underexplored in prior evaluations, and frames scientific inquiry as a multi-stage process involving database navigation, literature discovery under ambiguity, evidence consolidation, and structured synthesis.

% To operationalize these capabilities, we decompose scientific information seeking into four representative task types, each capturing a distinct stage of the scientific cognitive process with increasing complexity.
% The four task types are defined as follows.

\begin{comment}
\textsc{SciExplore} conceptualizes scientific information seeking as a progressive reasoning process that unfolds across multiple levels of abstraction.
Starting from the resolution of individual entities, we gradually increase the evaluation difficulty from document-level identification and evidence-level grounding to domain-level knowledge synthesis, and instantiate each capability through a corresponding task formulation.

Specifically, entity-level reasoning is examined through scientific database navigation, where agents must explore and traverse structured resources in the natural sciences.
Document-level identification and evidence-level grounding are evaluated via ambiguous literature retrieval and missing reference completion tasks, which require models to identify relevant papers under uncertainty and correctly align scientific claims with their supporting evidence.
Finally, domain-level synthesis is assessed through cross-source structured knowledge synthesis, where agents are required to integrate information from multiple heterogeneous studies and organize the resulting knowledge into a structured representation.
\end{comment}

\begin{comment}
As illustrated in Table \ref{tab:sciexplore-task-taxonomy}, this benchmark for evaluating science information seeking encompasses four progressively complex task types, each corresponding to different capability levels and atomic capabilities essential for effective information retrieval and synthesis in scientific research.

The first two tasks in our benchmark focus on database navigation and retrieval within scientific literature, respectively. Specifically, the Scientific Database Navigation task involves entity-level reasoning, assisting researchers to identify, filter, and disambiguate relevant entities while navigating through complex databases. This is crucial for researchers who must efficiently locate specific data amidst vast resources. Complementing this, the Ambiguous Literature Retrieval task targets document-level identification by helping users disambiguate intent and effectively discriminate among various documents to eliminate irrelevant ones. These tasks work together to evaluate the ability of LLMs or agents to find precise information in scientifically dense environments.

The third task, Missing Reference Completion, emphasizes evidence-level grounding. This task is essential for ensuring the integrity of claims made in scientific discourse. It involves processes such as claim abstraction, aligning claims with supporting evidence, and discovering and validating relevant evidence. This task not only aids in completing references but also enriches the overall quality of scientific research by ensuring that claims are adequately supported.

Finally, the Cross-source Structured Knowledge Synthesis task operates at the domain-level synthesis, representing a higher-order capability in our benchmark. This task involves domain abstraction and the integration of information from multiple sources to facilitate structured synthesis. It is critical for scholars who aim to produce comprehensive reviews or syntheses of existing literature, as it allows for the coherent amalgamation of knowledge across various studies, thus fostering a holistic understanding of specific scientific domains.
\end{comment}
% Table~\ref{tab:sciexplore-task-taxonomy} elaborates elaborates on the taxonomy of tasks in \textsc{SciExplore}.
% % As illustrated in Table~\ref{tab:sciexplore-task-taxonomy}, this benchmark for evaluating science information seeking includes four progressively complex task types, each corresponding to varying abilities essential for effective information retrieval and synthesis in scientific research.
% Specifically, T1, Scientific Database Navigation, helps researchers identify, filter, and disambiguate relevant entities in complex databases, which is crucial for efficiently locating specific data. T2, Ambiguous Literature Retrieval, aids users in disambiguating intent and discriminating among documents to eliminate irrelevant ones. Together, these tasks assess the ability of LLMs or agents to locate precise information in dense scientific environments.

% T3, Missing Reference Completion, ensures the integrity of scientific claims. This task includes claim abstraction, aligning claims with supporting evidence, and validating relevant evidence, thereby enriching the quality of research by ensuring adequate support for claims. Finally, T4, Cross-source Structured Knowledge Synthesis, represents a higher-order capability, involving domain abstraction, information extraction and the integration of information from multiple sources for structured synthesis. 


Table~\ref{tab:sciexplore-task-taxonomy} elaborates on the taxonomy of tasks in \textsc{SciExplore}, which spans a progression from entity-level retrieval to cross-source structured synthesis and is designed to jointly cover the core information-seeking competencies required in authentic scientific research workflows. 

Specifically, T1, Scientific Database Navigation, helps researchers identify, filter, and disambiguate relevant entities in complex databases, which is crucial for efficiently locating specific data. T2, Ambiguous Literature Retrieval, aids users in disambiguating intent and discriminating among documents to eliminate irrelevant ones. Together, these tasks assess the ability of LLMs or agents to locate precise information in dense scientific environments.
T3, Missing Reference Completion, ensures the integrity of scientific claims. This task includes claim abstraction, aligning claims with supporting evidence, and validating relevant evidence, thereby enriching the quality of research by ensuring adequate support for claims. Finally, T4, Cross-source Structured Knowledge Synthesis, represents a higher-order capability, involving domain abstraction, information extraction, and the integration of information from multiple sources for structured synthesis. Collectively, T3 and T4 probe the evidence-grounding and cross-source integration competencies that go beyond mere information location, complementing T1--T2 and together forming a comprehensive evaluation suite that mirrors the end-to-end workflow of a scientific research assistant.

% This task is critical for scholars aiming to produce comprehensive reviews or syntheses, as it facilitates the coherent amalgamation of knowledge from various studies, fostering a holistic understanding of scientific domains.


% Scientific information seeking in real research workflows is inherently progressive, requiring agents to advance from structured entity retrieval to document discovery, contextual evidence grounding, and ultimately domain-level knowledge synthesis.
% Guided by this observation, \textsc{SciExplore} is designed as a unified evaluation framework in which task definitions and data construction are jointly specified, so that each task instance instantiates a distinct level of scientific reasoning difficulty.
% All tasks are manually constructed by domain experts to ensure realism, answer uniqueness, and resistance to shortcut retrieval.

% \textsc{SciExplore} comprises four task types, ordered by increasing abstraction and integration demands, corresponding to successive stages of scientific information seeking. Details of the full dataset construction process are provided in the Appendix~\ref{app:datasets_construction}.

% \noindent{\textbf{Scientific Database Navigation.}}
% At the entity level, agents must recover uniquely identifiable targets through structured reasoning over interconnected scientific databases.
% To enforce genuine multi-hop traversal, task instances are constructed using a \emph{Reverse Trajectory Construction Strategy}, in which explicit entities are progressively masked and replaced with nested attribute-based constraints, such that correct resolution requires reversing the underlying dependency chain.

% \noindent{\textbf{Ambiguous Literature Retrieval.}}
% At the document level, agents are required to identify a specific target paper from imprecise and noisy descriptions, reflecting exploratory literature search under uncertainty.
% Construction adopts a two-stage design that combines \emph{Feature Denoising and Fuzzification} to substantially enlarge the retrieval space with \emph{Validation Constraint Injection} to preserve answer uniqueness through post-retrieval verification cues.

% \noindent{\textbf{Missing Citation Completion.}}
% At the argument level, evaluation shifts from document identification to contextual evidence grounding.
% Given a scientific paragraph with missing citations, agents must infer which prior works substantiate individual claims.
% To prevent shortcut matching, task contexts are rewritten to minimize surface-level overlap with the target references, ensuring that successful citation reconstruction depends on reasoning over claim--evidence relationships.

% \noindent{\textbf{Cross-source Structured Knowledge Synthesis.}}
% At the domain level, agents are challenged to integrate evidence across multiple studies into a structured representation.
% Given a research topic and a predefined comparison schema, tasks are constructed from expert-curated sources with strict format and schema constraints, such that accurate completion requires precise cross-paper evidence integration rather than partial or approximate extraction.

\begin{table*}[t]
\centering
\caption{Task taxonomy of \textsc{SciExplore}, detailing the task types, corresponding capability levels, atomic capabilities, task counts, and associated metrics.}
\label{tab:sciexplore-task-taxonomy}
% \vspace{-0.5em}
\scriptsize
% \footnotesize
\begin{tabular*}{\textwidth}{l @{\extracolsep{\fill}} l @{} l @{} l @{} c @{} l}
\toprule
\textbf{No.} & \textbf{Task Type} & \textbf{Capability Level} & \textbf{Atomic Capabilities} & \textbf{Tasks} & \textbf{Metric} \\
\midrule
T1 & \makecell[l]{Scientific Database \\ Navigation} & \makecell[l]{Entity-Level \\Reasoning} & \makecell[l]{Entity identification, filtering, disambiguation, \\ multi-constraint navigation} & 39 & Accuracy \\
\midrule
T2 & \makecell[l]{Ambiguous Literature \\ Retrieval} & \makecell[l]{Document-Level\\Identification} & \makecell[l]{Intent disambiguation, \\ document-level discrimination and elimination} & 32 & Accuracy \\
\midrule
T3 & \makecell[l]{Missing Reference \\Completion} & \makecell[l]{Evidence-Level\\Grounding} & \makecell[l]{Claim abstraction, claim--evidence alignment, \\ evidence discovery and validation} & 14 & Accuracy \\
\midrule
T4 & \makecell[l]{Cross-source Structured \\ Knowledge  Synthesis} & \makecell[l]{Domain-Level \\ Synthesis} & \makecell[l]{Domain abstraction, information extraction, \\ cross-source integration, structured synthesis} & 18 & Recall \\
\bottomrule
\end{tabular*}
% \vspace{-1em}
\end{table*}

\subsection{Dataset Construction}
\label{sec:dataset_construction}

\textsc{SciExplore} is constructed through a fully expert-driven process to ensure realism, answer uniqueness, and resistance to shortcut retrieval.
All task instances are manually curated by domain experts and designed to reflect authentic scientific research behaviors, rather than synthetic or template-based generation, as shown in Fig.~\ref{fig:dataset-pipeline}.

\vspace{0.5em}
\noindent{\textbf{Scientific Database Navigation.}} For T1, we design tasks that require genuine multi-step database navigation using a \emph{Reverse Trajectory Construction Strategy}, where explicit entity identifiers are progressively replaced by nested attribute-based constraints.
This design enforces structured reasoning over interconnected database records and prevents direct lookup.

\noindent{\textbf{Ambiguous Literature Retrieval.}} In T2, experts apply \emph{Feature Denoising and Fuzzification} to deliberately obscure surface-level cues, while preserving answer uniqueness through \emph{Validation Constraint Injection}, which enables post-retrieval verification without guiding the search process.

\noindent{\textbf{Missing Reference Completion.}} For T3 tasks, we construct citation reconstruction tasks that require aligning scientific claims with their appropriate supporting references.
Surface-level textual overlap is minimized, ensuring that successful prediction depends on understanding claim–evidence relationships rather than keyword matching.

\noindent{\textbf{Cross-Source Structured Knowledge Synthesis.} Tasks in T4 are derived from expert-curated comparison schema and require agents to integrate heterogeneous evidence across multiple papers into a strictly formatted table.
This setting emphasizes cross-source consistency, schema compliance, and high-precision synthesis.}

A detailed description of the construction procedure and validation criteria for each task is provided in Appendix~\ref{app:datasets_construction}.


% \paragraph{Multi-hop Database Navigation.}
% This task evaluates structured reasoning over interconnected scientific databases. Given a query that specifies a target entity indirectly through multiple attribute-based constraints, the agent must traverse several database entries to recover the unique target. Successful completion requires stepwise logical navigation rather than direct keyword lookup.

% \paragraph{Ambiguous Literature Retrieval.}
% This task models exploratory literature search under vague or noisy descriptions. The agent is given an imprecise methodological or conceptual description corresponding to a single target paper and must identify the correct work through broad search and post-retrieval verification. The output is a unique bibliographic record.

% \paragraph{Contextual Citation Reconstruction.}
% This task assesses an agent’s ability to infer supporting literature from scientific context. Given a paragraph with missing citations, the agent must analyze the claims and retrieve the appropriate references that substantiate each claim, relying on contextual reasoning rather than surface-level matching.

% \paragraph{Comparative Knowledge Synthesis.}
% This task represents domain-level synthesis. Given a research topic and a predefined comparison schema, the agent must identify relevant papers, extract heterogeneous methodological details, and synthesize them into a structured comparison table. Evaluation emphasizes both entity coverage and the accuracy of cross-source integration.

% \noindent
% Together, these tasks form a progressive evaluation framework, spanning from entity-level database navigation to document identification, contextual evidence reasoning, and domain-wide structured knowledge synthesis.


% \subsection{Dataset Construction}

% The construction of \textsc{SciExplore} follows a principled, expert-driven methodology designed to faithfully reflect the cognitive demands faced by autonomous scientific research assistants. All task instances are manually crafted by domain experts to ensure realism, difficulty, and alignment with authentic scientific workflows. Rather than relying on automatic generation, our construction process emphasizes controllability, answer uniqueness, and resistance to shortcut retrieval.

% \noindent \textbf{Multi-hop Database Navigation.}
% To construct tasks requiring multi-step structured reasoning, we adopt a \emph{Reverse Trajectory Construction Strategy}. Task generation begins from a uniquely identifiable target entity and progressively replaces explicit entity mentions with composite attribute-based descriptions extracted from scientific database records. During this process, related entities are recursively masked and substituted by their own attributes, forming nested constraints that encode deep dependency chains. Consequently, the final query can only be resolved by reversing the construction trajectory through stepwise database navigation, enforcing genuine multi-hop reasoning and preventing direct lookup.

% \noindent \textbf{Ambiguous Literature Retrieval.}
% Tasks involving ambiguous literature search are constructed to simulate early-stage exploratory research under uncertainty. Experts first extract core methodological or conceptual signals from a target paper and deliberately degrade them into vague, noisy descriptions through \emph{Feature Denoising and Fuzzification}, substantially enlarging the candidate retrieval space. To ensure answer uniqueness despite this ambiguity, we introduce \emph{Validation Constraint Injection}, adding auxiliary constraints that cannot directly guide retrieval but serve as strong post-retrieval validators (e.g., publication venue or author-level properties). This design compels agents to balance broad exploration with rigorous verification.

% \noindent \textbf{Contextual Citation Reconstruction.}
% For tasks requiring citation inference, we select citation-dense paragraphs from high-impact scientific literature and systematically rewrite the surrounding context to remove surface-level textual overlap with the original references. Overly canonical or survey-style citations are excluded to avoid trivial matching. This construction emphasizes reasoning over claim–evidence relationships, ensuring that correct reference reconstruction depends on contextual understanding rather than keyword similarity.

% \noindent \textbf{Comparative Knowledge Synthesis.}
% Tasks targeting structured synthesis are derived from expert-curated comparison tables in review papers across natural sciences and artificial intelligence. Given a predefined comparison schema, experts identify representative source papers and design tasks that require extracting heterogeneous methodological details scattered across multiple documents. Particular care is taken to include variations in terminology, implicit descriptions, and inconsistent reporting conventions. Strict schema and format constraints are enforced so that successful completion depends on accurate cross-source integration rather than partial or approximate extraction.

\begin{table*}[t]
\centering
\caption{Main results on the \textsc{SciExplore}. T4 is reported at two granularities: item- and row-level. All scores are percentages (\%).}
% \vspace{-0.5em}
\small
\label{tabs:main_results}

\begin{tabular*}{\textwidth}{l @{\extracolsep{\fill}} ccccc cc c}
\toprule

% \multirow{2}{*}{\textbf{Model}} & \multirow{2}{*}{\textbf{Category}} &
% \multirow{2}{*}{\textbf{T1}} &
% \multirow{2}{*}{\textbf{T2}} & \multirow{2}{*}{\textbf{T3}} &
% \multicolumn{2}{c}{\textbf{T4}} &
% \multirow{2}{*}{\textbf{Overall}} \\
\multirow{2}{*}[-0.8ex]{\textbf{Model}} & \multirow{2}{*}[-0.8ex]{\textbf{Category}} &
\multirow{2}{*}[-0.8ex]{\textbf{T1}} &
\multirow{2}{*}[-0.8ex]{\textbf{T2}} & \multirow{2}{*}[-0.8ex]{\textbf{T3}} &
\multicolumn{2}{c}{\textbf{T4}} &
\multirow{2}{*}[-0.8ex]{\textbf{Overall}} \\

\cmidrule(lr){6-7}
& & & & & {\scriptsize \textbf{Item}} & {\scriptsize \textbf{Row}} & \\
\midrule

\rowcolor{gray!15}
\multicolumn{8}{l}{\textit{Foundation LLMs}} \\

DeepSeek-V3.2
& Open-Source
& 20.51 & 6.25 & 5.95
& 8.32 & 2.25
& 11.06 \\

Qwen2.5-72B-Instruct
& Open-Source
& 7.69 & 0.00 & 0.00
& 8.63 & 6.94
& 4.99 \\

Qwen3-30B-A3B-Thinking
& Open-Source
& 10.26 & 0.00 & 3.57
& 2.54 & 1.60
& 4.50 \\

Qwen3-235B-A22B-Thinking
& Open-Source
& 20.51 & 3.12 & 0.00
& 7.85 & 4.15
& 7.87 \\

GPT-5.1
& Closed-Source
& 17.95 & 9.38 & 9.66
& 8.76 & 4.83
& 12.83 \\

Gemini-2.5-Pro
& Closed-Source
& 20.51 & 6.25 & 4.17
& 10.57 & 3.78
& 11.25 \\

Gemini-3-Pro
& Closed-Source
& 23.08 & 12.50 & 13.69
& 16.23 & 7.29
& 19.08 \\

\midrule
\rowcolor{gray!15}
\multicolumn{8}{l}{\textit{LLMs with Search Tools}} \\

DeepSeek-V3.2 w/ search
& Open-Source
& 17.95 & 9.38 & 9.66
& 9.20 & 3.95
& 11.12 \\

GPT-5.1 w/ search
& Closed-Source
& 15.38 & 34.38 & 20.85
& 12.13 & 5.12
& 21.61 \\

Gemini-3-Pro w/ search
& Closed-Source
& 33.33 & 53.12 & 49.40
& 23.53 & 11.49
& 46.46 \\

\midrule
\rowcolor{gray!15}
\multicolumn{8}{l}{\textit{Deep Research Agents}} \\

Tongyi-DeepResearch-30B-A3B
& Open-Source
& \textbf{43.59} & 50.00 & 34.40
& 11.12 & 4.26
& 36.93 \\

Gemini Deep Research
& Closed-Source
& 35.90 & \textbf{56.25} & 36.56
& 15.61 & 7.35
& 39.30 \\

OpenAI Deep Research
& Closed-Source
& 23.08 & 43.75 & \textbf{59.91}
& \textbf{30.14} & \textbf{18.59}
& \textbf{49.39} \\

\bottomrule
\end{tabular*}
% \vspace{-1em}
\end{table*}


\subsection{Quality Control}\label{sec:quality_control}

To ensure the rigor of \textsc{SciExplore}, we implemented a streamlined three-stage quality control process: query difficulty calibration, human cross-validation and answer uniqueness verification.

\noindent{\textbf{Query Difficulty Calibration.}}
To minimize parametric memory reliance, we strictly enforce two principles: \textit{Explicit Index Blocking} (removing unique identifiers to compel semantic matching) and \textit{Long-tail Entity Preference} (prioritizing obscure entities). This ensures tasks require deep retrieval reasoning rather than simple knowledge recall.

\noindent{\textbf{Human Cross-Validation.}} 
Annotators independently attempt to solve each question using search engines.
Questions are retained only if no solution can be found within a strict 10-minute limit, ensuring sufficient difficulty for the target models.

\noindent{\textbf{Answer Uniqueness Verification.}} 
We utilize search-enabled LLMs (e.g., Gemini-3-Pro, o3) to generate candidate solutions, which are then reviewed by human annotators. If any alternative answer is satisfied, the task is deemed ambiguous and discarded.

% After the filtering process, we obtain a high-quality set of 103 tasks, with 39, 32, 14, and 18 tasks in T1–T4, respectively.

\begin{comment}
\subsection{Source Statistics}
Unlike benchmarks grounded in a small set of homogeneous web sources, \textsc{SciExplore} spans over ten scientific sub-domains and integrates 16 authoritative databases with a pronounced long-tail distribution. This design forces models to navigate heterogeneous database structures and domain-specific metadata, thereby evaluating their robustness and generalization in authentic, interdisciplinary scientific workflows. Detailed statistics and distribution analyses are provided in Appendix~\ref{appendix_dataset_distribution}.

\subsection{Evaluation Metrics}

We employ task-specific evaluation metrics tailored to the structural characteristics of each task. T1 and T2 are evaluated using Exact Match accuracy based on semantic equivalence of the final answers. T3 is assessed with citation-level precision, measuring the proportion of correctly predicted references for each instance. For T4, we adopt a multi-granularity evaluation protocol at the row and item levels: the row-level metric measures recall over complete table rows, while the item-level metric evaluates fine-grained recall over individual cells or data items. 
An overall score is computed as a weighted sum of the four task scores to reflect holistic model performance. Details are provided in Appendix~\ref{app:evaluation_details}.
\end{comment}

\input{chaps/4_experiment}

\input{chaps/5_results}

\input{chaps/2_related_work}

\input{chaps/7_conclusion}

\bibliography{iclr2025_conference}
\bibliographystyle{iclr2025_conference}

\input{chaps/8_appendix}

\end{document}
